\section{Requirements}\label{sec:2}

\subsection{Functional Requirements, Acceptance Criteria and Use Cases}\label{sec:2.1}

\subsubsection{Functional Requirements}\label{sec:2.1.1}

{\small
\begin{longtable}[]{|>{\raggedright\arraybackslash}p{\dimexpr 0.0967\linewidth-2\tabcolsep\relax}|>{\raggedright\arraybackslash}p{\dimexpr 0.1827\linewidth-2\tabcolsep\relax}|>{\raggedright\arraybackslash}p{\dimexpr 0.6906\linewidth-2\tabcolsep\relax}|}
\hline
\textbf{ID} & \textbf{Group} & \textbf{Requirement} \\ \hline
\endhead
\textbf{FR-01} & Cattle \& Tag & Register an animal with its RFID tag
and maintain its farm/barn assignment and animal profile \\ \hline
\textbf{FR-02} & Cattle \& Tag & Record lost/damaged physical-tag
incidents while preserving the permanent RFID identity and complete
history. A replacement tag carries the same permanent identity;
assigning a different RFID is outside the scope of this phase. \\ \hline
\textbf{FR-03} & Cattle \& Tag & Deactivate a tag when an animal dies or
leaves the company \\ \hline
\textbf{FR-04} & Cattle \& Tag & Search animals by tag, ID, farm or
barn \\ \hline
\textbf{FR-05} & Gate \& Reader & Manage farms, barns and gates \\ \hline
\textbf{FR-06} & Gate \& Reader & Register readers and map each to a
gate and direction \\ \hline
\textbf{FR-07} & Gate \& Reader & Monitor reader status through
heartbeats and alert when a reader is offline \\ \hline
\textbf{FR-08} & Movement & Capture read events automatically (tag,
reader, timestamp) \\ \hline
\textbf{FR-09} & Movement & Collapse repeated observations of the same
physical passage using reader/gate context and passage/direction state.
Retain distinct passages even within the time window; deduplicate
transport retries by stable event ID. \\ \hline
\textbf{FR-10} & Movement & Determine the direction of each passage
(IN/OUT) \\ \hline
\textbf{FR-11} & Movement & Maintain a near real-time head count per
barn \\ \hline
\textbf{FR-12} & Movement & Record observed movement and approved
barn/farm transfers separately. Change assigned custody only through a
versioned transfer workflow; cross-farm transfer requires source and
destination approvals. \\ \hline
\textbf{FR-13} & Alerting & Raise a missing-animal alert when animals
that left have not returned by the end of the return window \\ \hline
\textbf{FR-14} & Alerting & Raise an unexpected-animal alert for
unregistered or deactivated tags \\ \hline
\textbf{FR-15} & Alerting & Raise a wrong-barn alert when an animal
enters a barn it is not assigned to \\ \hline
\textbf{FR-16} & Health & Record vaccinations (such as: vaccine, batch,
date, veterinarian) and vaccination schedules \\ \hline
\textbf{FR-17} & Health & Record health status, diagnoses and
treatments \\ \hline
\textbf{FR-18} & Health & View an animal' s complete
profile and history timeline \\ \hline
\textbf{FR-19} & Reporting & Generate head count, movement, vaccination
coverage, health and alert reports \\ \hline
\textbf{FR-20} & Reporting & Export authorized reports in CSV,
PDF and XLSX formats, preserving the same selected records and scope
filters. \\ \hline
\textbf{FR-21} & Reporting & Send relevant outbreak and
health-surveillance data to the National Animal Disease Control System
via API. \\ \hline
\textbf{FR-22} & Reporting & Provide an executive monitoring dashboard
across farms \\ \hline
\textbf{FR-23} & Auditing & Maintain an audit log of relevant system
activities \\ \hline
\textbf{FR-24} & Administration & Manage users, roles and permissions \\ \hline
\textbf{FR-25} & Administration & Configure approved return windows,
reader-offline timing, alert notification settings and
passage-observation windows; configuration must not discard genuine
passages or hide missing cattle. \\ \hline
\end{longtable}
}

\subsubsection{Acceptance Criteria and Use Cases}\label{sec:2.1.2}

Acceptance criteria define observable functional outcomes. The following
representative fixtures verify the functional groups; timing, faults,
permissions and evidence completeness are additionally verified by the
QAS and SC tables. All implementation tests are planned for M2.

{\small
\begin{longtable}[]{|>{\raggedright\arraybackslash}p{\dimexpr 0.1633\linewidth-2\tabcolsep\relax}|>{\raggedright\arraybackslash}p{\dimexpr 0.8067\linewidth-2\tabcolsep\relax}|}
\hline
\textbf{Category} & \textbf{Acceptance Criteria} \\ \hline
\endhead
\textbf{Cattle \& Tag} & Register a previously unused permanent RFID
once; reject an active collision. Search returns the same cattle/history
after an identity-preserving physical-tag incident or deactivation. No
permanent RFID is reassigned to another animal. \\ \hline
\textbf{Gate \& Reader} & Provision an approved farm/barn/gate/reader
binding; reject inconsistent bindings. A heartbeat timeout records a
reader-offline alert under the declared timing policy; an available
reader resumes capture. \\ \hline
\textbf{Movement} & From declared initial occupancy, one valid OUT
followed by IN restores the count. Replayed event IDs have one effect.
OUT then IN within the observation window retains both passages.
Approved transfers retain source/destination history; an unexpected gate
read cannot silently change custody. \\ \hline
\textbf{Alerting} & An outstanding OUT at return-window close creates a
scoped missing alert; OUT/IN/OUT is outstanding again. An unregistered
tag or wrong-barn entry creates its respective alert. Approved custody
transfers are excluded from expected-return obligations; stale context
yields provisional alerts. \\ \hline
\textbf{Health} & One authorized vaccination/treatment intent and three
retries create one record. A schedule retains its due date/status and
links a completed dose without duplicate vaccination effects.
Corrections append linked versions; complete permitted history remains
retrievable. \\ \hline
\textbf{Reporting} & Export the same authorized dataset to CSV/PDF/XLSX
with identical record/scope selection. Executive views exclude clinical
narratives. Final external acceptance marks completion; 202 remains
pending; ambiguous outcomes preserve the approved intent. \\ \hline
\textbf{Auditing} & Every covered committed operation has valid source
evidence and eventually canonical evidence after collector recovery.
User and machine actors are identifiable; denied access never causes
unauthorized effects. \\ \hline
\textbf{Admin\-is\-tration} & Unapproved role/configuration changes and
self-escalation are denied. FM/EM approvals and SA application are
linked to scope and policy version; revoked access fails at report
execution/download. \\ \hline
\end{longtable}
}

\emph{\Cref{app:A} describes five core use cases, including approved
transfer and external-submission subflows. \Cref{fig:1} shows external
human, reader and National-system participation; timers and Edge Core
are internal behavior. Participation never grants unrestricted
permission.}

\subsection{Non-Functional Requirements}\label{sec:2.2}

VinaBeef must process gate reads within two seconds, preserve movement
and health records during normal operation, and enforce RBAC with
comprehensive audit logging. These constraints guide the
architecture alongside the current scale of three farms, 120 barns and
approximately 25,000 cattle. API exchange is an explicit function;
continuity, changeability and exchange quality require additional design
criteria.

E identifies an explicit assignment requirement, I a justified inference
and P a project-defined criterion. Timing boundaries, permissions and test
budgets marked P are set by this design because the assignment does not
specify them; they are verified in M2 and kept separate from the
assignment' s mandatory outcomes.

{\small
\begin{longtable}[]{|>{\raggedright\arraybackslash}p{\dimexpr 0.2426\linewidth-2\tabcolsep\relax}|>{\raggedright\arraybackslash}p{\dimexpr 0.2425\linewidth-2\tabcolsep\relax}|>{\raggedright\arraybackslash}p{\dimexpr 0.2425\linewidth-2\tabcolsep\relax}|>{\raggedright\arraybackslash}p{\dimexpr 0.2425\linewidth-2\tabcolsep\relax}|}
\hline
\textbf{ID} & \textbf{Characteristic} & \textbf{Required outcome and scope} & \textbf{Origin} \\ \hline
\endhead
NFR-01 & Performance & Each required gate read completes correctly
within \(\leq\)2 seconds under the declared peak workload & E
deadline; P completion boundary \\ \hline
NFR-02 & Reliability and durability & Zero lost movement or health
records during normal operation; correct effects and preserved history &
E no-loss; I/P recovery extension \\ \hline
NFR-03 & Availability and continuity & Local gate processing continues
during a 60-minute WAN outage; its backlog recovers within 10
minutes while live traffic continues & I continuity; P budgets \\ \hline
NFR-04 & Security & Enforce the four human roles with action and
barn/farm scope; deny unauthorized access and changes & E RBAC; P
detailed policy \\ \hline
NFR-05 & Auditability & Every covered committed operation has valid
durable evidence; sensitive access is also traceable & E comprehensive
audit; P coverage measure \\ \hline
NFR-06 & Interoperability & Exchange approved data through APIs with
durable submission state, validated formats and safe handling of
uncertain outcomes & E API function; I/P quality \\ \hline
NFR-07 & Scalability and configurability & Support the current dataset
and declared gate workload; compatible farm/barn/reader additions use
configuration and provisioning & E scale; I/P expansion \\ \hline
NFR-08 & Maintainability and testability & Compatible reader/API adapter
changes preserve core business rules and pass relevant contract/domain
tests & E SOLID/M2; I/P change criteria \\ \hline
\end{longtable}
}

\subsubsection{Performance and Record Preservation}\label{sec:2.2.1}

The assignment' s approximately 200 cattle per farm in 30
minutes gives 200/1800 = 0.111 passages/s. This average does not
establish instantaneous raw-read load: one passage can produce repeated
observations, and cattle can arrive in bursts. M2 workload declarations
must include read multiplicity, burst timing, concurrent farms and
available resources.

The measurement starts at read creation at the source and ends
at the correct durable local outcome: movement, count effect and
required evidence, or a valid duplicate/invalid-input disposition. A
queued message alone does not complete processing. Every required read
must meet the two-second limit; percentiles are diagnostic. Missing,
incorrect and timed-out outcomes remain in the result, with clock
uncertainty accounted for.

Each persistent owner commits its business effect, operation identity
and required evidence atomically, then reports success or acknowledges
delivery. Stable identities and idempotency prevent replay from
duplicating movements, vaccinations or treatments. Health retries after
an uncertain response resolve the original intent, and corrections
retain linked versions. Separately stored projections need their own
transaction and recovery boundaries; broker confirmation does not prove
application completion.

The no-loss requirement covers both movement and health. Process
restart/replay and WAN outage are design extensions covered while durable
storage remains intact and correctly configured. Lossless capture during
full-disk, power or storage destruction requires additional protection;
unrestricted disaster zero-loss is not assumed.

\subsubsection{Availability and External Exchange}\label{sec:2.2.2}

The offline subset includes gate capture, local movement/count,
alert evaluation with sufficient cached context, and durable pending
synchronization. Stale or unknown context is marked provisional.
Clinical operations, transfers, configuration changes and new human
login remain online-only. Existing local dashboard reads expire within 15 minutes of the last central policy validation, or earlier
session expiry; machine capture uses a separate provisioned identity and
policy.

At one logical movement per passage, a 60-minute test at 1/9 movement/s
creates 400 movements per farm, or 1,200 across three farms. Draining
this backlog in 600 seconds while accepting live traffic requires at
least 400/600 + 1/9 = 7/9 movement/s per farm at a movement-processing
stage. This is a necessary lower bound, not proof of end-to-end
capacity. The outage/recovery budgets are finite test budgets, not
production SLAs.

External submissions are persisted before dispatch. The integration
adapter distinguishes pending, unknown and completed outcomes. HTTP 202
is pending acceptance; a timeout can follow a receiver side effect.
Ambiguous operations are retried automatically only when the receiver
contract establishes safety; otherwise they require status lookup or
reconciliation. Backoff respects applicable Retry-After, and bounded
retries retain unresolved intents. The actual National API
contract must be qualified during implementation.

\subsubsection{Security and the RBAC Matrix}\label{sec:2.2.3}

Security is the characteristic; RBAC is an authorization mechanism.
Chapter 4 lists Authorization among cross-cutting concerns and defines
Security in terms of appropriate authorized access.
The matrix implements the assignment' s four
roles. Unspecified actions are denied, and permission is checked against
the canonical resource at the service, query and export boundaries.

BS = Barn Supervisor; FM = Farm Manager; SA = System Administrator; EM =
Executive Management. B/F means assigned barn/farm scope; H-B/H-F
permits complete history of cattle currently in that scope; O-F permits
own-origin records after a transfer; C is company technical scope; A is
an approved aggregate view; D means denied. Each permission also
requires the applicable approval, version and workflow checks.

{\footnotesize
\begin{longtable}[]{|>{\raggedright\arraybackslash}p{\dimexpr 0.1940\linewidth-2\tabcolsep\relax}|>{\raggedright\arraybackslash}p{\dimexpr 0.1940\linewidth-2\tabcolsep\relax}|>{\raggedright\arraybackslash}p{\dimexpr 0.1940\linewidth-2\tabcolsep\relax}|>{\raggedright\arraybackslash}p{\dimexpr 0.1940\linewidth-2\tabcolsep\relax}|>{\raggedright\arraybackslash}p{\dimexpr 0.1940\linewidth-2\tabcolsep\relax}|}
\hline
\textbf{Action group} & \textbf{BS} & \textbf{FM} & \textbf{SA} & \textbf{EM} \\ \hline
\endhead
Read gate count and alert dashboard & B & F & D & A \\ \hline
Read cattle identity and movement history & H-B & H-F or O-F & D & A \\ \hline
Read full vaccination and health profile & H-B & H-F or O-F & D & D \\ \hline
Register cattle and permanent RFID & D & F & D & D \\ \hline
Record vaccinations, treatments or schedules & H-B & H-F & D & D \\ \hline
Request health correction & H-B & H-F or O-F & D & D \\ \hline
Approve and apply health correction & D & O-F & D & D \\ \hline
Append metadata or movement correction & D & F or O-F & D & D \\ \hline
Request transfer & B & F & D & D \\ \hline
Approve barn or farm transfer & D & F & D & D \\ \hline
Acknowledge or request alert resolution & B & F & D & D \\ \hline
Resolve alert or approve domain configuration & D & F & D & D \\ \hline
Apply approved technical configuration & D & D & C & D \\ \hline
Generate and download business reports & B & F & D & A \\ \hline
Approve external submission or reconciliation & D & F & D & D \\ \hline
Read protected business audit evidence & B & O-F & D & A \\ \hline
Manage accounts and role grants & D & Approve F grants & Apply approved
C actions & Approve company SA/EM grants \\ \hline
Hard delete committed movement, health or audit & D & D & D & D \\ \hline
\end{longtable}
}

Cross-farm transfers require both source- and destination-scope
approvals. New health writes require current custody, a stable intent
and the expected version. Health correction approval belongs to the
original record' s farm; ordinary corrections cannot
change custody. FM approves business configuration before SA applies it.
Role grants require independent FM/EM approval and prohibit
self-escalation. SA has scoped technical/security audit access, without
ordinary clinical read/write privilege. Report execution and download
recheck current scope; aggregate permissions exclude clinical
narratives.

Provisioned readers and processing/synchronization/integration workers
have separate machine identities limited to their assigned duties; they
cannot impersonate human approvals. Compatible farm onboarding uses
preapproved deployment provisioning. Runtime farm-creation actions
require a separately defined approval policy and remain denied while
unmapped.

\subsubsection{Audit Coverage}\label{sec:2.2.4}

Auditability expresses the project' s traceability goal.
It is supported by Chapter 4' s Accountability and
Nonrepudiation concerns; logging contributes evidence
but does not alone establish strong nonrepudiation. The following
coverage groups operationalize comprehensive audit logging.

{\small
\begin{longtable}[]{|>{\raggedright\arraybackslash}p{\dimexpr 0.3234\linewidth-2\tabcolsep\relax}|>{\raggedright\arraybackslash}p{\dimexpr 0.3233\linewidth-2\tabcolsep\relax}|>{\raggedright\arraybackslash}p{\dimexpr 0.3233\linewidth-2\tabcolsep\relax}|}
\hline
\textbf{Covered event group} & \textbf{Evidence to retain} & \textbf{Responsible source} \\ \hline
\endhead
Cattle registration and device/configuration changes & Identity/scope,
validated versions, actor and approvals & Registry and
device/configuration owner \\ \hline
Movement, count, alert and transfer lifecycle & Read/operation
references, direction, effects, transition reasons and approvals &
Movement, Count and Alert owners \\ \hline
Vaccination, health, treatment and correction & Stable intent, protected
record/version links, reason and actor; clinical-access evidence &
Health owner \\ \hline
Authentication, authorization, accounts and grants & Decision,
role/resource scope, policy version, approver and applier & IAM and
resource owner \\ \hline
Reports, exports and audit access & Authorized scope, job/artifact,
execution/download/access result & Reporting and Audit owners \\ \hline
External submission and reconciliation & Intent, destination/schema,
attempt, status and acceptance evidence & Integration owner \\ \hline
Synchronization, evidence controls and restore & Source/receipt
identity, recovery/configuration reference and operator/result & Sync,
Audit and storage owners \\ \hline
\end{longtable}
}

Common fields identify who or which device/service acted, what action
affected which resource, when it occurred, its outcome and its
operation/correlation identity. Clinical text, secrets and credentials
are excluded from broad technical logs; protected references preserve
provenance.

Required business evidence/outbox commits with the mutation at its
source. Remote collection can recover later; failure to persist
mandatory source evidence prevents an audited business success. Relay
duplicates retain stable audit identities and are deduplicated.
Denied access remains denied if logging fails, and any missing required
evidence is an audit gap. Application roles cannot overwrite/delete
canonical facts; protection against unrestricted infrastructure
administrators requires separate controls.
